\chapter{Fortran resources on the Web}

The following resources were available when this tutorial was revised (2026). Web addresses change, so if a link does not work, search for the name.

\section{The language}

\begin{itemize}
    \item \href{https://fortran-lang.org}{fortran-lang.org}: the home page of the Fortran community, with tutorials on modern Fortran, a list of compilers, and a \href{https://fortran-lang.org/packages/}{catalog of Fortran packages}.
    \item \href{https://fortranwiki.org}{Fortran Wiki}: an open wiki with articles on all aspects of Fortran, including Fortran 77.
    \item \href{https://fortran-lang.discourse.group}{Fortran Discourse}: a forum where you can ask questions about Fortran.
    \item \href{https://wg5-fortran.org/ARCHIVE/Fortran77.html}{The Fortran 77 standard}: the complete text of the standard (ANSI X3.9-1978), on the site of \href{https://wg5-fortran.org}{WG5}, the ISO committee for Fortran. The U.S. committee \href{https://j3-fortran.org}{J3} also publishes the drafts of the current standards.
    \item \href{https://www.nsc.liu.se/~boein/f77to90/f77to90.html}{Fortran 90 for the Fortran 77 Programmer}: a tutorial by Bo Einarsson and Yurij Shokin on the new features of Fortran 90, written for readers who already know Fortran 77.
\end{itemize}

\section{Compilers and tools}

\begin{itemize}
    \item \href{https://gcc.gnu.org/onlinedocs/gfortran/}{gfortran documentation}: the complete manual of gfortran, including all compiler options.
    \item \href{https://www.msys2.org}{MSYS2}: an easy way to install gfortran on Windows (see Chapter 4).
    \item Other free compilers: the \href{https://www.intel.com/content/www/us/en/developer/tools/oneapi/fortran-compiler.html}{Intel Fortran compiler} (\texttt{ifx}), \href{https://flang.llvm.org}{LLVM Flang}, and \href{https://lfortran.org}{LFortran}, which is still under development. Compiling your programs with more than one compiler is a good way to find errors and non-portable code.
    \item \href{https://www.gnu.org/software/gdb/}{gdb}, the GNU debugger, and \href{https://valgrind.org}{Valgrind} (see Chapter 20).
\end{itemize}

\section{Numerical software}

\begin{itemize}
    \item \href{https://www.netlib.org}{Netlib}, with BLAS, LAPACK and many other packages (Chapter 18).
    \item \href{https://www.netlib.org/lapack/lug/}{The LAPACK Users' Guide}.
    \item \href{https://gams.nist.gov}{GAMS}, the Guide to Available Mathematical Software.
    \item \href{https://www.openblas.net}{OpenBLAS}, a fast, free implementation of BLAS and LAPACK.
\end{itemize}

\section{The original tutorial}

This tutorial is based on the Fortran 77 tutorial written at Stanford University by Erik Boman, Sarah T. Whitlock and Paul H. Hargrove, which is still available at \url{https://web.stanford.edu/class/me200c/tutorial_77/}.
